\documentclass[10pt,a4paper]{amsart}
\usepackage[margin=19mm]{geometry}
\usepackage{amsmath,amssymb,mathtools}
\usepackage[scr=rsfs,cal=euler]{mathalfa}
\usepackage{xfrac}
\usepackage[unicode,hidelinks,bookmarks=false]{hyperref}
\newcommand{\ie}{i.e.\ }
\newcommand{\cf}{cf.\ }
\newcommand{\resp}{respectively\ }
\newcommand{\Subsection}{\subsection}
\providecommand{\nobreakdash}{\nobreak\hbox{-}}
\setlength{\emergencystretch}{2em}
\multlinegap0pt
\usepackage{amssymb}
\usepackage{xcolor}
\newtheorem{lemma}{Lemma}
\title{Fischer decompositions for entire functions of sufficiently low order}
\author{J.M. Aldaz and H. Render}
\date{Source: arXiv:2403.10419v3 (CC BY 4.0)}
\begin{document}
\maketitle

\noindent\emph{Source and licence.} Fischer decompositions for entire functions of sufficiently low order. Version arXiv:2403.10419v3 (CC BY 4.0), distributed by arXiv under the Creative Commons Attribution 4.0 International licence, http://creativecommons.org/licenses/by/4.0/, as recorded in the arXiv licence metadata for this exact version and shown on the abstract page https://arxiv.org/abs/2403.10419v3. Full licence notice: \texttt{licenses/arxiv-2403.10419-CC-BY-4.0.txt}.\par\smallskip

\noindent\textit{Editorial scope.} Counted unit: the article's lemma lemmaSeq (source0.tex lines 180--204). The article's complete original proof of it is delivered with it (source0.tex lines 206--391, one \texttt{proof} environment of 186 lines). No mathematics is written, repaired or re-derived: the statement and the proof are the article's own lines, in the article's own notation, and no retained line is substituted or re-flowed.  No further source-local context is needed: every cross-reference of every retained line resolves inside the two delivered spans.  Nothing was omitted from any retained span, so the package carries no omission record.  No retained line contains a citation, so the wrapper carries no bibliography and no bibliography entry.  No retained line contains a figure environment, an image-inclusion command or drawing code, so no image is delivered and the package declares no asset dependency.  Editorial formatting only, in addition: an AMS \texttt{amsart} preamble that restates the article's own declarations and macros used by the retained lines, namely the environments \texttt{lemma} on separate counters, together with the standard packages those lines need.  The retained environments print in this document as Lemma 1 in order of appearance, whereas the article prints the same items with the numbering of its own full document.  The complete native source of the article as distributed in the arXiv e-print for this exact version is bundled unchanged as \texttt{sources/156167/source0.tex}.
\begin{lemma}
	\label{lemmaSeq} Let $\left\{  a_{m}\right\}  _{m\in\mathbb{N}}$ be a sequence
	of non-negative numbers. Suppose there is a finite non-empty subset
	$E\subset\mathbb{N}\setminus\{0\}$ satisfying the following two conditions:
	
	(i) There exist constants $A \ge 1$,   $D \ge 0$ and $\alpha\in\mathbb{R}$ such that for
	all $m\in\mathbb{N}$,
	\[
	a_{m}\leq A\left(  m+D\right)  ^{\alpha}\max_{j\in E}a_{m+j}.
	\]
	(ii) There exist constants $A_{0},b_{0}>0$,  $D_{0},\alpha_{0}\geq0$ and
	$\sigma\neq0$ such that for all $m\in\mathbb{N}\setminus\{0\}$,
	\[
	a_{m}\leq A_{0}\ \frac{\left(  m+D_{0}\right)  ^{\alpha_{0}}}{b_{0}^{m}%
	}\ m^{-\frac{m}{\sigma}}\text{ }.
	\]
	Let $\beta_{\ast}:=\min E$ and $\beta^{\ast}=\max E.$ If
	\[
	0\leq\alpha<\frac{\beta_{\ast}}{\sigma},\text{ \quad or }\quad\alpha<0\text{
		and }\alpha<\frac{\beta^{\ast}}{\sigma},
	\]
	then for all $m\in\mathbb{N}$ we have $a_{m}=0$. When $\alpha
	=\frac{\beta_{\ast}}{\sigma} > 0$, the same conclusion holds under the
	additional hypothesis that $A < b_{0}^{\beta_{\ast}}.$
\end{lemma}

\begin{proof}
	Assume that $a_{m} > 0$ for some $m \ge 0$. We derive a contradiction. If $m =0$, by condition $(i)$ there exists a $j \in E$ such that $a_{m + j} > 0$, so we may suppose that $m > 0$. Since
	\begin{equation}
		a_{m}\leq A\left(  m+D\right)  ^{\alpha}\max_{j\in E}a_{m+j},\label{eq1}%
	\end{equation}
	there exists an $l_{1}\in E$ with
	\[
	a_{m}\leq A\left(  m+D\right)  ^{\alpha}a_{m+l_{1}}.
	\]
	Applying inequality (\ref{eq1}) to $a_{m+l_{1}}$ instead of $a_{m}$, we
	see that there exists an $l_{2}\in E$ such that
	\[
	a_{m}\leq A^{2}\left(  m+D\right)  ^{\alpha}\left(  m+l_{1}+D\right)
	^{\alpha}a_{m+l_{1}+l_{2}}.
	\]
	Now proceed inductively, to conclude that there exist $l_{1}, \dots,l_{j}\in E$
	satisfying
	\[
	a_{m}\leq A^{j}\left(  m+D\right)  ^{\alpha}\left(  m+l_{1}+D\right)
	^{\alpha}\cdots\left(  m+D+l_{1}+\cdots+l_{j-1}\right)  ^{\alpha}%
	a_{m+l_{1}+\cdots+l_{j}}.
	\]
	Define for $j \ge 1$, 
	\[
	p_{m,j} :=\left(  m+D\right)  \left(  m+l_{1}+D\right)  \cdots\left(
	m+D+l_{1}+\cdots+l_{j-1}\right),
	\]
	and $k_{j}:=
	m+l_{1}+\cdots+l_{j}.$ Our assumption $a_{m}\leq A_{0}\frac{\left(
		m+D_{0}\right)  ^{\alpha_{0}}}{b_{0}^{m}}\frac{1}{m^{\frac{m}{\sigma}}}$
	implies that
	\[
	a_{k_{j}}\leq A_{0}\frac{\left(  k_{j}+D_{0}\right)  ^{\alpha_{0}}}%
	{b_{0}^{k_{j}}}k_{j}^{-\frac{k_{j}}{\sigma}}, \
	\]
	so we have
	\[
	a_{m}\leq A_{0}\frac{\left(  k_{j}+D_{0}\right)  ^{\alpha_{0}}}{b_{0}^{k_{j}}%
	}A^{j}p_{m,j}^{\alpha}k_{j}^{-\frac{k_{j}}{\sigma}}.
	\]
	Taking $k_{j}$-th roots in the last inequality shows that
	\begin{equation}
		a_{m}^{\frac{1}{k_{j}}}\leq A_{0}^{\frac{1}{k_{j}}}\frac{\left(  k_{j}%
			+D_{0}\right)  ^{\frac{\alpha_{0}}{k_{j}}}}{b_{0}}A^{\frac{j}{k_{j}}}%
		p_{m,j}^{\frac{\alpha}{k_{j}}}k_{j}^{-\frac{1}{\sigma}}.\label{eqNEW}%
	\end{equation}
	Since $k_{j}=m+l_{1}+\cdots+l_{j}\geq m+\beta_{\ast}j > \beta_{\ast}j \ge j$,   we see
	that as $j \to \infty$, $k_{j}\rightarrow\infty$, so
	$a_{m}^{\frac{1}{k_{j}}}\rightarrow 1$,  $A_{0}
	^{\frac{1}{k_{j}}}\rightarrow 1$, $\left(  k_{j}+D_{0}\right)
	^{\frac{\alpha_{0}}{k_{j}}}\rightarrow1$, and $A^{\frac{j}{k_{j}}} \le
	A^{\frac{1}{\beta_{\ast}}}$. Thus, it suffices  to prove
	that $p_{m,j}^{\frac{\alpha}{k_{j}}}k_{j}^{-\frac{1}{\sigma}}\rightarrow0$ in order to
	obtain a contradiction, since then the left hand side of inequality (\ref{eqNEW}) converges to
	$1$ and the right hand side to $0.$ 
	
	If $
	0\leq\alpha
	<
	\beta_{\ast}/\sigma$, then  $m+D+l_{1}+\cdots+l_{s}\leq D+k_{j}$ for
	$s=0, \dots ,j-1$ and $p_{m,j}\leq\left(  D+k_{j}\right)  ^{j},$ so
	\[
	p_{m,j}^{\frac{\alpha}{k_{j}}}k_{j}^{-\frac{1}{\sigma}}
	\leq
	\left(  D+k_{j}\right)
	^{\alpha\frac{j}{k_{j}}}k_{j}^{-\frac{1}{\sigma}}\leq\left(  D+k_{j}\right)
	^{\frac{\alpha}{\beta_{\ast}}}k_{j}^{-\frac{1}{\sigma}}=\left(  \frac{D+k_{j}
	}{k_{j}}\right)  ^{\frac{\alpha}{\beta_{\ast}}}k_{j}^{\left(  \frac{\alpha
		}{\beta_{\ast}}-\frac{1}{\sigma}\right)  }.
	\]
	Since $\frac{\alpha}{\beta_{\ast}}-\frac{1}{\sigma}<0$ the term on the right
	hand side converges to $0$. If  $
	0 <\alpha
	=
	\beta_{\ast}/\sigma$, then
	\[
	1=\lim\sup_{j\rightarrow\infty}a_{m}^{\frac{1}{k_{j}}}
	\leq
	\lim\sup
	_{j\rightarrow\infty}A_{0}^{\frac{1}{k_{j}}}\frac{\left(  k_{j}+D_{0}\right)
		^{\frac{\alpha_{0}}{k_{j}}}}{b_{0}}A^{\frac{j}{k_{j}}}\left(  \frac{D+k_{j}
	}{k_{j}}\right)  ^{\frac{\alpha}{\beta_{\ast}}}\leq\frac{1}{b_{0}}A^{\frac
		{1}{\beta_{\ast}}} < 1.
	\] 
	This finishes the proof when $\alpha\geq0.$
	
	Let us discuss next the case $\alpha<0.$ Since $k_{j}\geq\beta_{\ast}j\ $for
	$j\geq1$, we have
	\[
	p_{m,j}
	=
	\left(  m+D\right)  \left(  m+l_{1}+D\right)  \cdots\left(
	m+D+l_{1}+\cdots+l_{j-1}\right)  
	\geq
	\beta_{\ast}^{j-1}\left(  j-1\right)  !.
	\]
	From $\alpha<0$ we conclude that
	\[
	p_{m,j}^{\frac{\alpha}{k_{j}}}\leq\left(  \beta_{\ast}\right)^{\alpha
		\frac{j-1}{k_{j}}}\left[  \left(  j-1\right)  !\right]  ^{\frac{\alpha}{k_{j}%
	}}.
	\]
	Recall that  $\beta^{\ast} :=\max\left\{  l:l\in E\right\}.$ Since
	$k_{j}=m+l_{1}+\cdots+l_{j}$, it follows that
	\[
	\beta_{\ast}j\leq k_{j}\leq m+\beta^{\ast}j.
	\]
	Thus $\left(  \beta_{\ast}\right)  ^{\alpha\frac{j-1}{k_{j}}}$ is bounded above.
	Furthermore, a version of Stirling's formula 
	with explicit bounds for $m \ge 0$ tells us that
	\[
	m!\geq\left(  2\pi m\right)  ^{\frac{1}{2}}\left(  \frac{m}{e}\right)
	^{m}e^{\frac{1}{12m+1}}.
	\]
	Using $\alpha<0$ and the preceding formula with $m=j-1 \ge 1$ we obtain
	\[
	\left[  \left(  j-1\right)  !\right]  ^{\frac{\alpha}{k_{j}}}\leq\left(
	2\pi\left(  j-1\right)  \right)  ^{\frac{1}{2}\frac{\alpha}{k_{j}}}\left(
	\frac{j-1}{e}\right)  ^{\alpha\frac{j-1}{k_{j}}}e^{\frac{\alpha}{k_j (12\left(
			j-1\right)  +1)}}.
	\]
	Note that
	\[
	\left(  \frac{1}{e}\right)  ^{\alpha\frac{j-1}{k_{j}}}=e^{-\alpha\frac
		{j-1}{k_{j}}}\leq e^{-\frac{\alpha}{\beta_{\ast}}},
	\]
	so
	\[
	\lim\sup_{j\rightarrow\infty}p_{m,j}^{\frac{\alpha}{k_{j}}}k_{j}^{-\frac{1}{\sigma}
	}
	\leq
	\lim\sup_{j\rightarrow\infty}  \left(  \beta_{\ast}\right)^{\alpha
		\frac{j-1}{k_{j}}} e^{-\frac{\alpha}{\beta_{\ast}}}\left(
	j-1\right)  ^{\alpha\frac{j-1}{k_{j}}}k_{j}^{-\frac{1}{\sigma}}
	\]
	Since $k_{j}\leq m+\beta^{\ast}j$ we have
	\[
	\frac{k_{j}}{j-1}
	\leq
	\frac{m}
	{j-1}+\beta^{\ast}+\frac{1}{j - 1}\beta^{\ast}.
	\]
	By assumption $\alpha  < \beta^{\ast}/  \sigma$, where $\alpha < 0$ (and we can have either $\sigma < 0$ or $\sigma > 0$). Let $\varepsilon>0$ be such that $\alpha <  (1 + \varepsilon)\beta^{\ast}/\sigma$. Now there exists a $j_{\varepsilon}$ such that for all
	$j\geq j_{\varepsilon}$,
	\[
	\frac{k_{j}}{j-1}\leq\left(  1+\varepsilon\right)  \beta^{\ast}.
	\]
	Then whenever $j\geq j_{\varepsilon}$,
	\[
	j-1
	\geq
	k_{j}\left(
	1+\varepsilon\right)^{-1}\left(  \beta^{\ast}\right)^{-1}.
	\]
	Using again that $\alpha<0$ we get
	\begin{align*}
		\left(  j-1\right)  ^{\alpha\frac{j-1}{k_{j}}} &  \leq k_{j}^{\alpha\frac
			{j-1}{k_{j}}}\left(  1+\varepsilon\right)  ^{-\alpha\frac{j-1}{k_{j}}}\left(
		\beta^{\ast}\right)  ^{-\alpha\frac{j-1}{k_{j}}}\\
		&  \leq k_{j}^{\alpha\frac{j-1}{k_{j}}}\left(  1+\varepsilon\right)
		^{-\alpha\frac{1}{\beta_{\ast}}}\left(  \beta^{\ast}\right)  ^{-\alpha\frac
			{1}{\beta_{\ast}}}.
	\end{align*}
	Now
	\[
	\lim\sup_{j\rightarrow\infty}p_{m,j}^{\frac{\alpha}{k_{j}}}k_{j}^{-\frac{1}{\sigma}
	}
	\leq 
	e^{-\frac{\alpha}{\beta_{\ast}}}\left(  1+\varepsilon\right)
	^{-\alpha\frac{1}{\beta_{\ast}}}\left(  \beta^{\ast}\right)  ^{-\alpha\frac
		{1}{\beta_{\ast}}}
	\lim\sup_{j\rightarrow\infty}
	\left(  \beta_{\ast}\right)^{\alpha
		\frac{j-1}{k_{j}}}
	k_{j}^{\alpha\frac{j-1}{k_{j}}
	}k_{j}^{-\frac{1}{\sigma}},
	\]
	and this limit equals $0$ since
	\[
	\alpha\frac{j-1}{k_{j}}
	\leq
	\frac{\alpha}{\left(
		1+\varepsilon\right)  \beta^{\ast}}
	< \frac{1}{\sigma}.
	\]
	\end{proof}

\end{document}
